\subsection{Differences with existing algorithms}

\begin{frame}[t,fragile]{Execution policy parameter}
\begin{itemize}

  \item The execution policy parameter can be used with:
     \begin{itemize}
        \item Iterator based algorithms.
        \item Range based algorithms.
    \end{itemize}
    
  \mode<presentation>{\vfill\pause}
  \item New iterator based parallel algorithms with updated interfaces.

\begin{lstlisting}
ranges::unary_transform_result<I, O>
ranges::transform(Ep&& exec, I first, S last, O result, OutS result_last,
                  F op, Proj proj = {});
\end{lstlisting}

  \begin{itemize}
    \item \cppcode{Ep}: Execution policy.
    \item \cppcode{I}, \cppcode{S}: Input iterator and sentinel.
    \item \cppcode{O}, \cppcode{OutS}: Output iterator and sentinel.
    \item \cppcode{F}: Unary operation.
    \item \cppcode{Proj}: Projection (optional).
  \end{itemize}
  
\end{itemize}
\end{frame}

\begin{frame}[t,fragile]{Execution policy and ranges}
\begin{itemize}
    \item New range based parallel algorithms with updated interfaces.

\begin{lstlisting}
ranges::unary_transform_result<borrowed_iterator_t<R>, borrowed_iterator_t<OutR>>
ranges::transform(Ep&& exec, R&& r, OutR&& result_r, F op, Proj proj = {});
\end{lstlisting}

  \begin{itemize}
    \item \cppcode{Ep}: Execution policy.
    \item \cppcode{R}: Input range.
    \item \cppcode{OutR}: Output range.
    \item \cppcode{F}: Unary operation.
    \item \cppcode{Proj}: Projection (optional).
  \end{itemize}
\end{itemize}
\end{frame}

\begin{frame}[t,fragile]{Parallel algorithms require random access}
\begin{itemize}
    \item Iterator based parallel algorithms require random access iterators.
      \begin{itemize}
        \item Both input and output iterators must be random access.
      \end{itemize}

    \item Range based parallel algorithms require random access ranges.
      \begin{itemize}
        \item Both input and output ranges must be random access.
      \end{itemize}

    \mode<presentation>{\vfill\pause}
    \item \textmark{Consequences}:
      \begin{itemize}
        \item Cannot be used with input, forward, or bidirectional iterators/ranges.
        \item Cannot be used with output iterators/ranges that are not random access.
        \item Cannot be used with some views (e.g., \cppcode{std::views::filter}).
      \end{itemize}

    \mode<presentation>{\vfill\pause}
    \item \textmark{Why?}:
      \begin{itemize}
        \item Fast division of work into chunks.
        \item Random access allows for efficient partitioning and load balancing.
      \end{itemize}
\end{itemize}
\end{frame}

\begin{frame}[t,fragile]{Parallel algorithms require bounded ranges}
\begin{itemize}
  \item Iterator based parallel algorithms require sized sentinels.
    \begin{itemize}
      \item The distance between the input iterator and sentinel must be computable in constant time.
      \item The distance between the output iterator and sentinel must be computable in constant time.
    \end{itemize}

  \item Range based parallel algorithms require sized ranges.
    \begin{itemize}
      \item The size of the input range must be computable in constant time.
      \item The size of the output range must be computable in constant time.
    \end{itemize}

  \mode<presentation>{\vfill\pause}
  \item \textmark{Consequences}:
    \begin{itemize}
      \item Cannot be used with ranges that do not have a known size.
      \item Cannot be used with some views (e.g., \cppcode{std::views::filter}).
      \item When different sizes are involved processing stops when the end of the shorter range is reached.
    \end{itemize}

  \mode<presentation>{\vfill\pause}
  \item \textmark{Why?}:
    \begin{itemize}
        \item Efficiently divide work into chunks.
        \item Fixes safety problems with out of range access.
    \end{itemize}
\end{itemize}
\end{frame}

\begin{frame}[t,fragile]{Fixing a pending issue: output ranges}
\begin{itemize}
  \item Output is specified by a pair of iterators or a range.
\end{itemize}
  
\begin{lstlisting}
ranges::unary_transform_result<I, O>
ranges::transform(Ep&& exec, I first, S last, O result, @OutS result_last@,
                  F op, Proj proj = {});

ranges::unary_transform_result<borrowed_iterator_t<R>, borrowed_iterator_t<OutR>>
ranges::transform(Ep&& exec, R&& r, @OutR&& result_r@, F op, Proj proj = {});
\end{lstlisting}
\end{frame}